\chapter{Congenital Disorders}

\section*{Patient focus}
Congenital abnormalities are conditions present from birth. Many are detected before birth on pregnancy ultrasound---this allows parents and doctors to plan for treatment immediately after delivery. Some abnormalities are found on newborn screening or later in childhood. Imaging helps define the exact anatomy, assess severity, and plan surgery. Different scans are used depending on the body system involved: echocardiography for heart defects, ultrasound for abdominal problems, and MRI for brain and spine abnormalities.


\label{ch:ch68_congenital}

\section{Imaging Role}
Congenital disease requires age-appropriate anatomic and functional assessment, often across multiple modalities. Imaging may define anatomy, physiology, complications, and suitability for intervention.

\section{Prenatal Imaging}

\subsection{Ultrasound Screening}
\begin{itemize}
  \item \textbf{First trimester} (11--13 weeks): nuchal translucency, nasal bone, major structural anomalies
  \item \textbf{Second trimester} (18--22 weeks): detailed anatomical survey---brain, spine, heart, kidneys, limbs, placenta, amniotic fluid
  \item \textbf{Third trimester}: growth assessment, follow-up of detected anomalies, placental position
\end{itemize}

\subsection{Common Prenatal Ultrasound Findings}
\begin{table}[h]
\centering
\caption{Common congenital anomalies on prenatal US}
\begin{tabular}{|l|l|l|}
\hline
\textbf{System} & \textbf{Anomaly} & \textbf{US findings} \\
\hline
Cardiac & CHD (various) & Abnormal 4-chamber view, outflow tract anomalies \\
CNS & Neural tube defect & Spina bifida, anencephaly, encephalocele \\
CNS & Ventriculomegaly & Lateral ventricle $>$10 mm \\
GI & Gastroschisis & Bowel outside abdominal wall, no membrane \\
GI & Omphalocele & Bowel/liver in sac at umbilicus \\
Renal & Pelviectasis & Anteroposterior diameter $>$10 mm \\
Skeletal & Limb reduction & Absent or short limbs \\
\hline
\end{tabular}
\end{table}

\subsection{Fetal MRI Indications}
Fetal MRI is recommended when ultrasound findings are equivocal or when detailed brain/spine assessment is needed:
\begin{enumerate}
  \item \textbf{Brain anomalies}: ventriculomegaly, agenesis of corpus callosum, cortical malformations, posterior fossa abnormalities (Dandy-Walker, Chiari)
  \item \textbf{Spine}: suspected spinal dysraphism, sacrococcygeal teratoma
  \item \textbf{Chest}: congenital diaphragmatic hernia---assessing lung volume and liver position
  \item \textbf{Body wall}: complex gastroschisis, assessment of bowel viability
\end{enumerate}

\begin{pearlbox}[Fetal MRI safety]
Fetal MRI uses fast sequences (SSFP, single-shot T2) and is considered safe in the second and third trimesters. No gadolinium is given to the mother. MRI provides superior soft tissue contrast compared to ultrasound, particularly for brain and spine abnormalities, and helps guide parental counselling and delivery planning.
\end{pearlbox}

\section{Cardiac Congenital Abnormalities}

\begin{table}[h]
\centering
\caption{Common congenital heart defects}
\begin{tabular}{|l|l|l|}
\hline
\textbf{Defect} & \textbf{Key imaging findings} & \textbf{Modality} \\
\hline
ASD & Defect in interatrial septum, RA$\rightarrow$LA shunt & Echo \\
Coarctation of aorta & Focal narrowing at isthmus, collateral vessels & Echo, CT, MRA \\
HLHS & Hypoplastic left heart, mitral/aortic atresia & Echo \\
Tetralogy of Fallot & VSD, overriding aorta, RVOT obstruction, RV hypertrophy & Echo, CT \\
Transposition of great arteries & Aorta anterior from RV, PA posterior from LV & Echo, CT \\
VSD & Defect in interventricular septum, LV$\rightarrow$RV shunt & Echo \\
\hline
\end{tabular}
\end{table}

\textbf{Postnatal evaluation}: Echocardiography is the primary investigation for suspected CHD. CT angiography provides excellent spatial resolution for complex anatomy and surgical planning. MRI provides functional assessment (flow quantification, ventricular volumes).

\section{Gastrointestinal Congenital Abnormalities}

\subsection{Biliary Atresia}
\begin{itemize}
  \item \textbf{US}: absent or small gallbladder, triangular cord sign (echogenic tissue around portal vein), absent gallbladder-portal vein swing sign
  \item \textbf{HIDA scan}: no excretion of tracer into bowel---but not specific
  \item \textbf{Diagnosis}: operative cholangiography is gold standard; early diagnosis (before 60 days) improves Kasai procedure outcomes
\end{itemize}

\subsection{Pyloric Stenosis}
\begin{itemize}
  \item \textbf{US}: pyloric muscle thickness $>$4 mm, pyloric channel length $>$16 mm, ``cervix sign'' on longitudinal view
  \item \textbf{Timing}: typically presents at 2--8 weeks with non-bilious projectile vomiting
  \item \textbf{Treatment}: pyloromyotomy---US diagnosis avoids barium study and its aspiration risk
\end{itemize}

\subsection{Malrotation}
\begin{itemize}
  \item \textbf{Prenatal US}: abnormal vessel relationship (SMA/SMV reversal), dilated bowel loops
  \item \textbf{Postnatal}: upper GI series is gold standard---duodenaljejunal junction should be to the left of spine at the level of duodenal bulb
  \item \textbf{CT}: mesenteric swirling, reversed SMA/SMV relationship
\end{itemize}

\section{Renal Congenital Abnormalities}

\begin{table}[h]
\centering
\caption{Common renal congenital anomalies}
\begin{tabular}{|l|l|l|}
\hline
\textbf{Anomaly} & \textbf{US/CT findings} & \textbf{Clinical significance} \\
\hline
Duplex system & Two collecting systems on one side & May have ureterocele, VUR \\
Horseshoe kidney & Fusion of lower poles, anterior to aorta & Usually asymptomatic; higher stone risk \\
Multicystic dysplastic kidney & Multiple non-communicating cysts, no function & Contralateral kidney usually normal \\
Pelvic kidney & Ectopic kidney in pelvis & Usually asymptomatic \\
\hline
\end{tabular}
\end{table}

\section{Skeletal Dysplasias}
\begin{itemize}
  \item \textbf{Skeletal survey}: full-body radiographs (AP and lateral) using low-dose technique
  \item \textbf{Common dysplasias}: achondroplasia (short limbs, narrow pedicles, trident hand), osteogenesis imperfecta (multiple fractures, bowing, osteopenia), thanatophoric dysplasia (curved femurs, ``telephone receiver'' appearance)
  \item \textbf{3D CT reconstruction}: helpful for surgical planning in complex skeletal anomalies
\end{itemize}

\section{CNS Congenital Abnormalities}

\begin{table}[h]
\centering
\caption{Common CNS congenital anomalies}
\begin{tabular}{|l|l|l|}
\hline
\textbf{Anomaly} & \textbf{Imaging findings} & \textbf{Modality} \\
\hline
Agenesis of corpus callosum & Absent CC, parallel ventricles & MRI \\
Dandy-Walker & Enlarged posterior fossa, cyst, vermian hypoplasia & US, MRI \\
Holoprosencephaly & Failure of forebrain division & US, MRI \\
Lissencephaly & Smooth brain, thick cortex & MRI \\
Myelomeningocele & Spinal dysraphism, Chiari II & US (prenatal), MRI \\
\hline
\end{tabular}
\end{table}

\section{When to Recommend Postnatal Imaging}
\begin{enumerate}
  \item Antenatal US anomaly detected $\rightarrow$ postnatal confirmatory imaging (US, echo, MRI) within first 48 hours
  \item Suspected CHD (cyanosis, murmur, abnormal pulse oximetry) $\rightarrow$ urgent echocardiography
  \item Bilious vomiting in neonate $\rightarrow$ abdominal X-ray (``double bubble'' in duodenal atresia), upper GI series (malrotation/volvulus)
  \item Spina bifida $\rightarrow$ spinal US (infants $<$3 months) or MRI (older children)
  \item Craniofacial anomalies $\rightarrow$ CT for bony detail, MRI for brain assessment
\end{enumerate}

\section{Interpretation}
Use segmental analysis, laterality, connections, flow, obstruction, shunts, and associated anomalies. Compare with prior studies and state whether findings are new, expected, or postoperative.

\section{Safety}
Minimize radiation and sedation burden, use pediatric expertise, and consider pregnancy in adolescents. Acute cyanosis, obstruction, heart failure, vascular compromise, or airway risk requires urgent multidisciplinary communication.

\begin{pearlbox}[Congenital imaging safety]
Paediatric imaging requires dose reduction (ALARA principle), appropriate scanner selection (US preferred in neonates), and minimal sedation. MRI under general anaesthesia should be avoided when possible---fast MRI sequences may allow scanning without sedation. Always use age-appropriate protocols and sizes for equipment.
\end{pearlbox}

\clearpage
